\input{header.tex}
\input{newcommands.tex}

% Ausfüllbare Leistungskontrolle: Felder über \feld, \zeichen, \kreuz, \antwort (newcommands.tex).
% Rote Lehrerfelder über \lehrerpunkte und \lehrerfeld: für die Klasse gesperrt, das Skript
% scripts/lk_auswerten.py gibt sie frei und trägt Punkte ein (Lösungsdatei: loesung_beispiel.toml).
% Optionsfelder (radio) brauchen unicode=false.
\hypersetup{unicode=false, pdftitle={Leistungskontrolle}}

\begin{document}
\begin{Form}

Name: \feld[6cm]{name} \hfill Klasse: \feld[2cm]{klasse} \hfill Datum: \feld[3cm]{datum}

\begin{center}\textbf{\Large{Leistungskontrolle Thema}}\end{center}

{\small\itshape Klicke in die blauen Felder und schreibe oder kreuze dort deine Antworten an.
Die roten Felder füllt deine Lehrkraft aus.
Speichere die Datei vor dem Hochladen und öffne sie zur Kontrolle noch einmal.}

\bigskip
\begin{enumerate}[(1)]

	\item Kreuze an, welche Werte vom Datentyp \texttt{int} sind. \lehrerpunkte{L-a1}{2}\\
	\kreuz{a1-1}\ \texttt{15} \qquad \kreuz{a1-2}\ \texttt{1.5} \qquad
	\kreuz{a1-3}\ \texttt{"7"} \qquad \kreuz{a1-4}\ \texttt{-3}

	\item Entscheide für jede Aussage, ob sie richtig oder falsch ist. \lehrerpunkte{L-a2}{2}

	\begin{tabular}{p{10cm}cc}
		\toprule
		Aussage & richtig & falsch \\
		\midrule
		Eine Zeile mit \texttt{\#} ist ein Kommentar. & \kreuz{a2-1-r} & \kreuz{a2-1-f} \\
		\texttt{input()} liefert immer eine Zahl. & \kreuz{a2-2-r} & \kreuz{a2-2-f} \\
		\bottomrule
	\end{tabular}

	\item Ordne zu: Eingabe (E), Verarbeitung (V) oder Ausgabe (A). \lehrerpunkte{L-a3}{2}\\
	Münzen einwerfen \zeichen{a3-1} \qquad Wechselgeld berechnen \zeichen{a3-2}

	\item Gib an, was Python ausgibt. \lehrerpunkte{L-a4}{2}\\
	\texttt{17 // 5} = \feld[1.5cm]{a4-a} \qquad \texttt{17 \% 5} = \feld[1.5cm]{a4-b}

	\item Erkläre den Unterschied zwischen \texttt{=} und \texttt{==}. \lehrerpunkte{L-a5}{2}
	\antwort{a5}

\end{enumerate}

% Bewertung: füllt die Lehrkraft in der Korrekturfassung aus
\bigskip
\noindent\textcolor{akzentrot}{\rule{\linewidth}{0.6pt}}\\[0.2cm]
{\color{akzentrot}\textbf{Bewertung} \quad Punkte: \lehrerfeld[1.3cm]{L-summe}\,/\,10
\qquad Note: \lehrerfeld[1.3cm]{L-note} \qquad Bemerkung: \lehrerfeld[4.5cm]{L-bemerkung}}

% Lösungen NICHT in diese Datei, sondern in eine eigene Datei für die Lehrkraft
% (Erwartungsbild als PDF, Lösungsdatei für das Skript als .toml).

\end{Form}
\end{document}
